% celestial_moonshine_bridge.tex
%
% Celestial-Moonshine Bridge: the $w_{1+\infty}$ celestial chiral algebra of
% 4d self-dual gravity coupled to a K3 target recovers Mathieu moonshine in
% its matter sector. A chapter-quality inscription establishing three
% theorems (celestial-moonshine bridge; Mathieu--celestial correspondence;
% shadow-tower moonshine) plus one corollary promoting Vol III Mathieu
% moonshine to a theorem of the programme's celestial framework.
%
% Macro hygiene: every non-core macro guarded by \providecommand to survive
% standalone extraction and cross-volume migration.

\providecommand{\cA}{\mathcal A}
\providecommand{\cB}{\mathcal B}
\providecommand{\cC}{\mathcal C}
\providecommand{\cE}{\mathcal E}
\providecommand{\cF}{\mathcal F}
\providecommand{\cH}{\mathcal H}
\providecommand{\cI}{\mathcal I}
\providecommand{\cL}{\mathcal L}
\providecommand{\cM}{\mathcal M}
\providecommand{\cO}{\mathcal O}
\providecommand{\cP}{\mathcal P}
\providecommand{\cS}{\mathcal S}
\providecommand{\cT}{\mathcal T}
\providecommand{\cW}{\mathcal W}
\providecommand{\frakg}{\mathfrak g}
\providecommand{\frakk}{\mathfrak k}
\providecommand{\frakm}{\mathfrak m}
\providecommand{\RR}{\mathbb R}
\providecommand{\CC}{\mathbb C}
\providecommand{\PP}{\mathbb P}
\providecommand{\ZZ}{\mathbb Z}
\providecommand{\NN}{\mathbb N}
\providecommand{\HH}{\mathbb H}
\providecommand{\scri}{\mathscr I}
\providecommand{\Ran}{\operatorname{Ran}}
\providecommand{\Sym}{\operatorname{Sym}}
\providecommand{\Obs}{\operatorname{Obs}}
\providecommand{\End}{\operatorname{End}}
\providecommand{\Hom}{\operatorname{Hom}}
\providecommand{\Res}{\operatorname{Res}}
\providecommand{\Tr}{\operatorname{Tr}}
\providecommand{\Mel}{\widetilde}
\providecommand{\SCchtop}{\mathsf{SC}^{\mathrm{ch,top}}}
\providecommand{\Barch}{\mathrm B^{\mathrm{ch}}}
\providecommand{\BarchOrd}{\mathrm B^{\mathrm{ch},\mathrm{ord}}}
\providecommand{\ChirHoch}{\mathrm{ChirHoch}}
\providecommand{\dlog}{\mathrm d\log}
\providecommand{\delbar}{\bar\partial}
\providecommand{\wtc}{\mathfrak w_{1+\infty}}
\providecommand{\kapch}{\kappaChHodge}
\providecommand{\kapcat}{\kappa_{\mathrm{cat}}}
\providecommand{\kapBKM}{\kappa_{\mathrm{BKM}}}
\providecommand{\kapfiber}{\kappa_{\mathrm{fiber}}}
\providecommand{\cAcel}{\cA^{\mathrm{cel}}}
\providecommand{\MtwFour}{M_{24}}
\providecommand{\Zder}{Z^{\mathrm{der}}_{\mathrm{ch}}}
\providecommand{\ClaimStatusConditional}{\textnormal{[conditional]}}

%=======================================================================
\section{Celestial--moonshine bridge}
\label{sec:celestial-moonshine-bridge}
%=======================================================================

\noindent
On one side, celestial holography identifies the celestial chiral algebra of $4$d
self-dual gravity as $\wtc$, the universal higher-spin tower of soft
gravitons (Strominger~\cite{StromingerWInfinityCelestial},
Guevara--Himwich--Pate--Strominger~\cite{GuevaraHimwichPateStrominger21},
Bittleston--Heuveline--Skinner~\cite{BittlestonHeuvelineSkinner},
Costello--Paquette--Sharma~\cite{CostelloPaquetteSharmaCelestial}). On the
other, Mathieu moonshine identifies the K3 elliptic genus as a virtual character
of a graded $\MtwFour$-module
(Eguchi--Ooguri--Tachikawa~\cite{EguchiOoguriTachikawa10},
proved by Gannon~\cite{Gannon16}), with umbral extensions to the $23$
Niemeier lattices proved by Duncan--Griffin--Ono~\cite{DuncanGriffinOno15}
after the Cheng--Duncan--Harvey construction~\cite{ChengDuncanHarvey1,
ChengDuncanHarvey2}.

The two sides agree by construction once
the celestial dictionary of
Theorem~\ref{thm:uch-main} is applied to $4$d twistor gravity coupled
to a K3 target. The resulting celestial chiral algebra splits
canonically as an OPE tensor product
$\cAcel = \wtc \otimes \cA^{\mathcal{N}{=}4}_{K3}$.
For a fixed K3 sigma model, only its actual symmetry group
$G_\Sigma$ acts on the matter factor; Gaberdiel--Hohenegger--Volpato
\cite{GaberdielHoheneggerVolpato11} prove that no single K3 sigma
model realizes all of $\MtwFour$.  The full $\MtwFour$ statement used
below is therefore the virtual Mathieu module supplied by Gannon, not
a geometric action on one chiral algebra.

Every arrow is supplied by a theorem of this programme or cited from the literature
with a precise conversion identity; analogy carries no mathematical
weight.

\S\ref{subsec:cmb-setup} fixes the combined 4d twistor gravity + K3
target setup and records the tensor decomposition of the celestial
chiral algebra.
\S\ref{subsec:cmb-main} states the celestial--moonshine bridge theorem
and proves it from the universal celestial holography theorem.
\S\ref{subsec:cmb-twining} sharpens the bridge to an explicit
$G_\Sigma$-twined correspondence and to the virtual $\MtwFour$ twining
genera of Gannon.
\S\ref{subsec:cmb-shadow} identifies the shadow-tower coefficients of
the K3 chiral algebra with the polar and shadow data of the Mathieu
mock modular forms, and records the Rademacher comparison as a
separate hypothesis.
\S\ref{subsec:cmb-umbral} extends conditionally to the 23 Niemeier
lattices using the Duncan--Griffin--Ono theorem and an explicit
celestial realization datum.
\S\ref{subsec:cmb-corollary} records the in-programme derivation of Vol~III
Theorem~\ref{thm:mathieu-moonshine} from the celestial framework.

% =====================================
\subsection{Setup: 4d twistor gravity coupled to K3}
\label{subsec:cmb-setup}
% =====================================

Fix a self-dual twistor gravity theory $T_{\mathrm{sd grav}}$ in the
sense of Costello--Paquette--Sharma~\cite{CostelloPaquetteSharmaCelestial}
and Bittleston--Heuveline--Skinner~\cite{BittlestonHeuvelineSkinner},
coupled to a K3 nonlinear sigma model of target volume $1$ and
B-field class zero. The combined $4$d theory
\begin{equation}
\label{eq:cmb-combined-theory}
T_4^{K3}
\;:=\;
T_{\mathrm{sd grav}} \;\boxtimes\; \Sigma^{K3},
\end{equation}
with $\Sigma^{K3}$ the K3 worldsheet sigma model at the self-dual
radius, admits a holomorphic--topological twist $T_4^{K3,\mathrm{HT}}$
by Costello--Li~\cite{CL16}; the gravity factor is
HT-twistable by the Bittleston--Heuveline--Skinner / Costello--Paquette--Sharma
twistor construction, and the sigma-model factor is HT-twistable by
Costello--Gwilliam chiral quantization~\cite{CostelloGwilliamVolI,
CostelloGwilliamVolII} at the $(0,4)$ locus (Witten's A-twist on K3
produces a chiral N=4 at $c=6$).

\begin{definition}[Combined celestial chiral algebra]
\label{def:cmb-combined-cel-algebra}
Let $\cAcel(T_4^{K3}) := \cAcel(T_{\mathrm{sd grav}} \boxtimes \Sigma^{K3})$
be the celestial chiral algebra of the combined 4d theory in the sense
of Definition~\ref{def:uch-boundary-chiral-algebra}. By the tensor
structure of HT-twisted boundary algebras
(Costello--Gwilliam~\cite[\S 4]{CostelloGwilliamVolII}) and the
celestial-boundary product structure of
Theorem~\ref{thm:uch-main}(i), the combined celestial algebra splits
canonically as the OPE tensor product
\begin{equation}
\label{eq:cmb-tensor-decomposition}
\cAcel(T_4^{K3})
\;=\;
\cAcel(T_{\mathrm{sd grav}}) \;\otimes\; \cAcel(\Sigma^{K3})
\;=\;
\wtc \;\otimes\; \cA^{\mathcal{N}{=}4}_{K3},
\end{equation}
where $\wtc$ is the celestial $w_{1+\infty}$-algebra of Strominger
(Proposition~\ref{prop:uch-gravity-leading} applied to
$T_{\mathrm{sd grav}}$) and $\cA^{\mathcal{N}{=}4}_{K3}$ is the chiral
$\mathcal{N}{=}4$ superconformal algebra at $c = 6$ that arises as the
Costello--Gaiotto boundary algebra of the HT twist of the K3 sigma
model.
\end{definition}

\begin{remark}[Two sectors, two $\kapch$]
\label{rem:cmb-two-kappa}
The two factors carry distinct modular characteristics, subscripted throughout:
\[
\kapch(\wtc) \;=\; \wtc\text{'s chiral central charge }
 = c(\wtc) / 2 \qquad
\kapch(\cA^{\mathcal{N}{=}4}_{K3}) \;=\; \chi(\cO_{K3}) \;=\; 2.
\]
For the $\mathcal{N}{=}4$ K3 algebra the chiral-Hodge characteristic is
the arithmetic genus $\chi(\cO_{K3}) = 2$, the $\Phi$-functor value
$\kapch(A_{K3}) = \Xi(K3) = 2$, not the bare Virasoro central charge
$c/2 = 3$; it is this value that the polar coefficient of the K3 mock
modular form records, $h|_{q^{-1/8}} = -\kapch = -2$.
Every occurrence of $\kappa$ is subscripted: $\kapch$ for chiral bar,
$\kapcat$ for the categorified Euler
characteristic of K3 (which equals $\chi(\cO_{K3}) = 2$), and
$\kapBKM$ for the Borcherds lift sector of the attached BKM
algebra in the umbral extension of
\S\ref{subsec:cmb-umbral}.
\end{remark}

\begin{remark}[$\wtc$ carries no K3 moonshine action]
\label{rem:cmb-gravity-m24-invariant}
$T_{\mathrm{sd grav}}$ is pure gravity; its celestial algebra $\wtc$
does not carry the K3 symmetry action. If a chosen K3 sigma model
$\Sigma$ has symmetry group $G_\Sigma$, then $G_\Sigma$ acts trivially
on $\wtc$ and through its genuine sigma-model action on
\(\cA^{\mathcal{N}{=}4}_{K3,\Sigma}\). The full group $\MtwFour$ is
not a symmetry group of a single K3 sigma model; the $\MtwFour$ data
used in Mathieu moonshine are Gannon's virtual graded module and its
twining characters. Thus the celestial bridge has two honest forms:
geometric twining for \(g\in G_\Sigma\), and virtual $\MtwFour$
twining after the Gannon module is supplied as external input.
\end{remark}

\begin{remark}[Ordering convention]
\label{rem:cmb-no-ordering}
Every bar complex in this chapter is the symmetric (factorization) bar,
not the ordered bar: the spectral-parameter structure on $\PP^1_{\mathrm{cel}}$
is chiral-holomorphic, and no $E_1$
ordered bar enters. Where the word ``ordered'' appears in a cross-reference
it refers to labelled configurations on the celestial sphere,
not to a total order on $\RR$.
\end{remark}

% =====================================
\subsection{The bridge theorem}
\label{subsec:cmb-main}
% =====================================

Strominger's celestial $\wtc$ and the K3 matter chiral algebra assemble
into a tensor-product celestial algebra. Twining is geometric for the
actual symmetry group of the chosen K3 sigma model and virtual for the
full Mathieu group.

\begin{theorem}[Celestial--moonshine bridge; \ClaimStatusConditional]
\label{thm:celestial-moonshine-bridge}
\index{celestial holography!Mathieu moonshine}
\index{Mathieu moonshine!celestial origin}
Let $T_4^{K3} = T_{\mathrm{sd grav}} \boxtimes \Sigma^{K3}$ be the 4d
twistor gravity coupled to a K3 sigma model
(\S\ref{subsec:cmb-setup}). Let \(G_\Sigma\) be the actual symmetry group
of this sigma model, and let \(M=\bigoplus_{n\geq0}M_n\) denote the
graded virtual \(\MtwFour\)-module whose character is Gannon's Mathieu
moonshine theorem. There is a natural map of chiral algebras
\begin{equation}
\label{eq:cmb-bridge-map}
\mu
\;\colon\;
\cA^{\mathcal{N}{=}4}_{K3}
\;\hookrightarrow\;
\cAcel(T_4^{K3})
\;=\;
\wtc \otimes \cA^{\mathcal{N}{=}4}_{K3},
\qquad
\mu(a) = 1_{\wtc} \otimes a,
\end{equation}
induced by the Costello--Paquette celestial dictionary applied to the
boundary of the combined HT twist. The map $\mu$ satisfies:
\begin{enumerate}[label=\textup{(\roman*)}]
\item \emph{(Celestial correlator.)}
The celestial correlator of $n$ conformal primaries of $T_4^{K3}$ on
$\PP^1_{\mathrm{cel}}$ splits as
\[
\bigl\langle \textstyle\prod_{i=1}^n \cO_i
\bigr\rangle_{\cAcel(T_4^{K3})}
\;=\;
\bigl\langle \textstyle\prod_{i=1}^n \cO_i^{\mathrm{grav}}
\bigr\rangle_{\wtc}
\;\cdot\;
\bigl\langle \textstyle\prod_{i=1}^n \cO_i^{\mathrm{matter}}
\bigr\rangle_{\cA^{\mathcal{N}{=}4}_{K3}},
\]
i.e.\ the celestial partition function of $T_4^{K3}$ factors as
$Z_{\wtc}(\tau) \cdot Z_{K3}(\tau, z)$.
\item \emph{(Mathieu moonshine from twining.)}
For \(g\in G_\Sigma\), geometric twining acts trivially on the first
factor and by the K3 sigma-model action on the second:
\[
Z_{\cAcel}^{(g)}(\tau, z)
\;=\;
Z_{\wtc}(\tau) \cdot Z_{K3}^{(g)}(\tau, z).
\]
For the full group \(\MtwFour\), the same formula is a virtual
character identity after replacing the geometric K3 action by the
twining character \(\Tr_M(g\,q^{L_0-c/24}y^{J_0})\) supplied by
Gannon. It is not a statement that \(\MtwFour\) acts on one K3 sigma
model.
\item \emph{($\wtc$-module decomposition.)}
The tensor-product structure induces a decomposition of
$\cAcel(T_4^{K3})$ as a $\wtc \otimes \cA^{\mathcal{N}{=}4}_{K3}$-module;
restricted to the $\wtc$-action, it decomposes into copies of the
gravity module labelled by the K3 matter state space. Twined traces are
indexed by actual \(G_\Sigma\)-classes in the geometric case and by
Gannon's virtual \(\MtwFour\)-characters in the Mathieu case. The Frame
shape is a useful label for eta-products, but it is not an injective
label of \(\MtwFour\)-conjugacy classes.
\end{enumerate}
\end{theorem}

\begin{proof}
\emph{(i).}
By Theorem~\ref{thm:uch-main}(iii) the $n$-point celestial correlator
equals the $n$-point chiral factorization homology of the celestial
algebra on $\PP^1_{\mathrm{cel}}$. Chiral factorization homology is
multiplicative under OPE tensor product of chiral algebras
(Beilinson--Drinfeld~\cite[\S 3.4]{BeilinsonDrinfeldChiral}), and the
combined celestial algebra factors as $\wtc \otimes
\cA^{\mathcal{N}{=}4}_{K3}$ by
Definition~\ref{def:cmb-combined-cel-algebra}. The splitting of
correlators is the multiplicativity of factorization homology under the
tensor product.

\emph{(ii).}
For \(g\in G_\Sigma\), functoriality of the celestial dictionary in
\(Q\)-equivariant maps gives an action on the matter factor and the
trivial action on the gravity factor; hence the trace factors. GHV
classifies the symmetry groups of individual K3 sigma models and does
not supply an action of all of \(\MtwFour\) on one model. For the full
\(\MtwFour\) statement, the input is instead Gannon's virtual module;
the displayed identity is then a character identity, not a geometric
sigma-model symmetry statement.

\emph{(iii).}
The $\wtc$-module structure on $\cAcel(T_4^{K3})$ arises from the
first factor of the tensor product. A generic element of
$\cAcel(T_4^{K3})$ has the form $\sum_i w_i \otimes a_i$ with $w_i \in
\wtc$ and $a_i \in \cA^{\mathcal{N}{=}4}_{K3}$; as a $\wtc$-module it
is a direct sum of copies of $\wtc$ labelled by a basis of
$\cA^{\mathcal{N}{=}4}_{K3}$. Geometric twining acts only on the second
factor, so it changes the trace on the multiplicity space and leaves
the $\wtc$-module itself fixed. The same sentence holds for the virtual
Mathieu module after that module is supplied externally. Since distinct
conjugacy classes can have the same Frame shape, the label is the
twining character, not the Frame shape alone.
\end{proof}

\begin{remark}[Sources of each arrow]
\label{rem:cmb-construction-not-narration}
Every arrow in Theorem~\ref{thm:celestial-moonshine-bridge} is a theorem
of this programme
(Theorem~\ref{thm:uch-main}, Definition~\ref{def:cmb-combined-cel-algebra})
or a cited external result (Costello--Gwilliam chiral quantization of
HT twists; GHV classification of K3 sigma-model symmetries; Gannon's
Mathieu moonshine theorem; Beilinson--Drinfeld multiplicativity of
factorization homology). No arrow uses an all-\(\MtwFour\) action on a
single K3 sigma model.
\end{remark}

% =====================================
\subsection{Mathieu--celestial correspondence and explicit action on $\wtc$-modules}
\label{subsec:cmb-twining}
% =====================================

Theorem~\ref{thm:celestial-moonshine-bridge}(iii) decomposes
$\cAcel(T_4^{K3})$ into $\wtc$-module summands with K3 matter
multiplicity. The twined trace takes the following explicit form.

\begin{theorem}[Mathieu--celestial correspondence; \ClaimStatusConditional]
\label{thm:mathieu-celestial-correspondence}
\index{Mathieu moonshine!explicit $\wtc$-modules}
\index{celestial holography!$\MtwFour$-action}
Let $T_4^{K3}$ be as in \S\ref{subsec:cmb-setup}. For \(g\in G_\Sigma\),
or for \(g\in \MtwFour\) after the Gannon virtual module has been
supplied, let $\cM_g$ denote the corresponding twined $\wtc$-module
trace in Theorem~\ref{thm:celestial-moonshine-bridge}(iii).
Then:
\begin{enumerate}[label=\textup{(\roman*)}]
\item \emph{(Module character.)}
The graded character of $\cM_g$ as a $\wtc$-module equals the
K3 twining genus $Z_{K3}^{(g)}(\tau, z)$ times the $\wtc$-vacuum
character $Z_{\wtc}(\tau)$:
\[
\mathrm{ch}_{\cM_g}(\tau, z)
\;=\;
Z_{\wtc}(\tau) \cdot Z_{K3}^{(g)}(\tau, z).
\]
\item \emph{(Class $1A$ and the full elliptic genus.)}
For $g = 1A$ the Frame shape is $1^{24}$, $Z_{K3}^{(1A)} = Z_{K3}$ is
the full K3 elliptic genus, and $\cM_{1A}$ is the $\wtc$-module whose
character equals the product $Z_{\wtc}(\tau) \cdot Z_{K3}(\tau, z)$.
Evaluating at the Witten index point $z = 0$ recovers
$Z_{\wtc}(\tau) \cdot 24 = Z_{\wtc}(\tau) \cdot \chi(K3)$.
\item \emph{(Higher classes and Mathieu multiplicities.)}
For every Mathieu class \(g\), the $\mathcal{N}{=}4$ decomposition of
Gannon's virtual twining genus \(Z_{K3}^{(g)}\) into massless and massive
characters yields coefficients \(A_n^{(g)}=\Tr_{\rho_n}(g)\). If
\(g\in G_\Sigma\), this is also a geometric sigma-model twining trace.
\item \emph{(Frame-shape scope.)}
The assignment by Frame shape is not injective on \(\MtwFour\)-conjugacy
classes: the pairs \(7A/7B\), \(14A/14B\), \(15A/15B\), \(21A/21B\),
and \(23A/23B\) have coincident Frame-shape data. Thus the honest label
for the celestial trace is the twining character, not the Frame shape
alone. Injectivity is asserted only after passing to distinct twining
genera or after restricting to a class set where the relevant characters
are separated.
\end{enumerate}
\end{theorem}

\begin{proof}
\emph{(i).}
By Theorem~\ref{thm:celestial-moonshine-bridge}(i) the twined celestial
partition function factors as $Z_{\wtc}(\tau) \cdot Z_{K3}^{(g)}(\tau,
z)$; this equals the graded character of $\cM_g$ by the definition of
$\cM_g$ as the $\wtc$-module trace at label $g$ in the
Theorem~\ref{thm:celestial-moonshine-bridge}(iii) decomposition.

\emph{(ii).}
For $g = 1A$, the Frame shape $1^{24}$ is the identity label, so
$Z_{K3}^{(1A)} = Z_{K3} = 2\,\phi_{0,1}$
(Vol~III Definition~\ref{V3-def:k3-elliptic-genus}). Evaluating at
$z = 0$: $Z_{K3}(\tau, 0) = 24 = \chi(K3)$, and the Witten index of
$\cM_{1A}$ picks up the multiplication by $\chi(K3) = 24$.

\emph{(iii).}
The $\mathcal{N}{=}4$ decomposition of \(Z_{K3}^{(g)}\) is the Mathieu
moonshine decomposition
(Vol~III equation~\eqref{V3-eq:mathieu-decomposition}) twined by \(g\).
By Gannon~\cite{Gannon16}, the coefficients
\(A_n^{(g)}=\Tr_{\rho_n}(g)\) are character values of the virtual
\(\MtwFour\)-module. When \(g\) is geometrically realized at the chosen
K3 model, the same trace is an honest sigma-model twining trace.

\emph{(iv).}
The listed class pairs have the same Frame-shape data, so no argument
using \(a_1(\pi_g)\) can distinguish them.  The celestial construction
therefore records the full twined character \(Z_{K3}^{(g)}\).  Any
injectivity statement must be made at that character level or on a
subcollection where the characters are already known to be distinct.
\end{proof}

\begin{remark}[Twining acts on multiplicity spaces]
\label{rem:cmb-permutation-not-deformation}
Twining in Theorem~\ref{thm:mathieu-celestial-correspondence} changes
only the K3 multiplicity trace.  The $\wtc$-algebra is fixed pointwise,
consistent with Remark~\ref{rem:cmb-gravity-m24-invariant}.  The
Mathieu representation structure lives in the supplied matter
character, separate from any deformation of the stress tensor or of
the $\wtc$-algebra itself.
\end{remark}

% =====================================
\subsection{Shadow-tower moonshine: polar data and Rademacher comparison}
\label{subsec:cmb-shadow}
% =====================================

Volume~I records the shadow-invariant / $L$-function connection as
aspirational for the class M lane. Through the celestial--moonshine
bridge, the K3 matter factor supplies the polar and Zwegers-shadow data
of the Mathieu mock modular forms. Identifying the full tower of
bar-shadow coefficients with Rademacher coefficients is an additional
comparison problem, not a consequence of the tensor decomposition alone.

\begin{theorem}[Shadow-tower moonshine comparison; \ClaimStatusConditional]
\label{thm:shadow-tower-moonshine}
\index{shadow tower!Mathieu moonshine}
\index{mock modular form!shadow tower}
Let $\cA^{\mathcal{N}{=}4}_{K3}$ be the celestial K3 chiral algebra of
\S\ref{subsec:cmb-setup}. Let \(h_g(\tau)\) be the twined Mathieu mock
modular form attached to \(g\) by Gannon's virtual module. Then:
\begin{enumerate}[label=\textup{(\roman*)}]
\item The polar term of every \(h_g\) at \(q^{-1/8}\) has coefficient
\(-2\). It is the trace of \(g\) on the trivial two-dimensional polar
piece of the Mathieu module, not \(-\Tr_{\widetilde H(K3,\ZZ)}(g)/12\).
\item The Zwegers shadow of \(h_g\) is proportional to
\(\Tr_{\mathbf{24}}(g)\eta(\tau)^3\). It is not the Frame-shape eta
product \(\eta_g=\prod_i\eta(i\tau)^{a_i}\), which belongs to the
umbral eta-product package.
\item A comparison map
\[
\mathcal R_n^{\mathrm{K3}}\colon
S_n^{(g)}(\cA^{\mathcal{N}{=}4}_{K3})
\longrightarrow [q^n]\,h_g(\tau)
\]
from bar-shadow depth to the \(n\)-th Rademacher/Fourier coefficient is
not constructed here. When such maps are supplied and are compatible
with the polar and shadow data above, the celestial K3 lane recovers
the Mathieu coefficient \([q^n]h_g\). Without \(\mathcal R_n^{K3}\)
there is no theorem identifying bar depth \(r\) with a Rademacher
``depth''.
\end{enumerate}
\end{theorem}

\begin{proof}
The Mathieu module has a two-dimensional polar piece on which every
\(\MtwFour\)-element acts trivially; hence the polar coefficient is
\(-2\) for all \(g\). The shadow formula is the standard Mathieu
moonshine shadow formula: the twining parameter enters through the
trace of \(g\) on the \(24\)-dimensional permutation representation,
giving \(\Tr_{\mathbf{24}}(g)\eta^3\). Frame-shape eta products have
different weights and occur in the umbral eta-product package, so they
cannot be identified with this Zwegers shadow. Rademacher series
reconstruct Fourier coefficients from polar and shadow data, but that
reconstruction is indexed by Fourier exponent \(n\), not by chiral bar
depth \(r\). The final clause is therefore exactly the additional
comparison datum needed to turn the polar/shadow agreement into a
bar-depth theorem.
\end{proof}

\begin{remark}[Scope inherited from the UCH machinery]
\label{rem:shadow-tower-moonshine-uch-scope}
The theorem above uses only the polar and shadow data of the Mathieu
mock modular forms. The stronger old assertion
\(S_r^{(g)}=\mathrm{Rad}_r[h_g]\) is not stated: neither
\(\mathrm{Rad}_r\) nor a canonical map from chiral bar depth to
Rademacher Fourier index has been constructed in this chapter. A strict
chain-level reading at class-\(\mathbf M\) remains tied to the
weight-completed class-\(\mathbf M\) comparison theorem
\(\ref{thm:uch-gravity-chain-level}\) and to the additional maps
\(\mathcal R_n^{\mathrm{K3}}\).
\end{remark}

\begin{remark}[Shadow/$L$-function closure in the K3 lane]
\label{rem:cmb-shadow-L-function-closure}
Volume~I observed that the shadow $L$-function of Vir and of a
generic class~M chiral algebra has trivial zeros at negative integers;
the connection to \(F_g\leftrightarrow L^{\mathrm{sh}}(1-2g)\) was
flagged as an aspiration. The K3 lane supplies a precise polar/shadow
comparison with the Mathieu mock modular forms. The full identification
of Rademacher coefficients with bar-depth coefficients remains
conditional on the maps \(\mathcal R_n^{\mathrm{K3}}\); the general
class~M case remains open.
\end{remark}

% =====================================
\subsection{Umbral extension: the 23 Niemeier lattices}
\label{subsec:cmb-umbral}
% =====================================

Cheng--Duncan--Harvey~\cite{ChengDuncanHarvey1,ChengDuncanHarvey2}
constructed the umbral moonshine package attached to the $23$ Niemeier
lattices with nonempty root system, and Duncan--Griffin--Ono
\cite{DuncanGriffinOno15} proved the umbral moonshine conjecture. For
each such Niemeier lattice $N$ there is an umbral group $G_N$, a
lambency \(\ell(N)\), and a mock modular form whose coefficients are
$G_N$-representation dimensions. The \(A_1^{24}\) case recovers Mathieu
moonshine.

\begin{proposition}[Umbral celestial correspondence; \ClaimStatusConditional]
\label{prop:umbral-celestial-correspondence}
\index{Umbral moonshine!celestial origin}
Let $N$ be a Niemeier lattice with nonempty root system, $G_N$ its
umbral group. Assume an umbral celestial realization datum
\[
\mathfrak U_N=(\Sigma^N,\cA_N^{\mathcal N=4},G_N\curvearrowright
\cA_N^{\mathcal N=4},\mathcal R^N_n)
\]
whose matter chiral algebra carries the \(G_N\)-module structure of
the Duncan--Griffin--Ono theorem and whose comparison maps
\(\mathcal R^N_n\) identify the matter twined traces with the umbral
mock modular forms. Then the combined 4d theory
$T_{\mathrm{sd grav}} \boxtimes \Sigma^N$ has celestial chiral algebra
\[
\cAcel(T_{\mathrm{sd grav}} \boxtimes \Sigma^N)
\;=\;
\wtc \otimes \cA^{\mathcal{N}{=}4}_{N},
\]
with $\cA^{\mathcal{N}{=}4}_{N}$ the chiral $\mathcal{N}{=}4$ algebra at
$c = 6$ supplied by \(\mathfrak U_N\). The \(G_N\)-twined celestial
correlators then reproduce the umbral moonshine coefficients.
\end{proposition}

\begin{proof}
The construction is identical to
Definition~\ref{def:cmb-combined-cel-algebra} with $\Sigma^{K3}$
replaced by $\Sigma^N$. The \(G_N\)-action and the comparison with the
umbral mock modular forms are not consequences of GHV K3 sigma-model
symmetry. They are exactly the data included in \(\mathfrak U_N\),
whose representation-theoretic part is the Duncan--Griffin--Ono
theorem. The tensor decomposition, factorization of correlators, and
twining factorization \(Z_{G_N}^{(g)}=Z_{\wtc}\cdot Z_N^{(g)}\) then
follow from the same proof as
Theorem~\ref{thm:celestial-moonshine-bridge}.
\end{proof}

\begin{remark}[Six routes to umbral moonshine, one celestial output]
\label{rem:cmb-six-routes-one-output}
Volume~III records six constructions of $G(K3 \times E)$;
their convergence is the content of CY-C, conjectural. The
celestial correspondence is a seventh route,
with output a celestial chiral algebra (not a BKM algebra).
Convergence of the celestial output with the six Vol~III outputs is
a conjecture about the celestial dictionary composed with the CY-to-chiral
functor $\Phi$. The present statement is that celestial holography
applied to $T_{\mathrm{sd grav}} \boxtimes \Sigma^N$
produces a $\wtc \otimes \cA^{\mathcal{N}{=}4}_{N}$ celestial algebra
carrying a $G_N$-action; identification with any of the six Vol~III
algebraizations is a separate question.
\end{remark}

\begin{remark}[Reach of the umbral extension]
\label{rem:cmb-umbral-reach}
Proposition~\ref{prop:umbral-celestial-correspondence} covers the $23$
Niemeier lattices with nonempty root system. The $24$th Niemeier
lattice (Leech lattice $\Lambda_{24}$, no roots) lies outside the
umbral moonshine framework; its celestial analogue is Conway
moonshine~\cite{DuncanMackOno} on the Conway module $V^{s\natural}$ at
$c = 12$, distinct from K3-type celestial algebras at $c = 6$.
Extending the celestial--moonshine bridge to $c = 12$ uses a
different 4d HT theory ($T_{\mathrm{sd grav}} \boxtimes V^{s\natural}$
as a chiral algebra coupled to gravity), beyond
the present chapter.
\end{remark}

% =====================================
\subsection{Mathieu moonshine as a conditional celestial trace statement}
\label{subsec:cmb-corollary}
% =====================================

\begin{corollary}[Celestial placement of Mathieu twining characters;
\ClaimStatusConditional]
\label{cor:celestial-holography-recovers-moonshine}
\index{Mathieu moonshine!celestial holography corollary}
The universal celestial holography functor
$\Phi_{\mathrm{hol}} \colon T_4 \mapsto \cAcel(T_4)$
(Theorem~\ref{thm:uch-main}), applied to the K3-coupled 4d twistor
gravity $T_4^{K3} = T_{\mathrm{sd grav}} \boxtimes \Sigma^{K3}$,
outputs a celestial chiral algebra
$\cAcel(T_4^{K3}) = \wtc \otimes \cA^{\mathcal{N}{=}4}_{K3}$.  For
$g \in G_\Sigma$, the geometric twined trace is
$Z_{\wtc}\cdot Z_{K3,\Sigma}^{(g)}$.  For $g\in \MtwFour$, after
supplying Gannon's virtual Mathieu module, the same tensor-product
trace reads the Mathieu twining character as a K3 matter trace.  The
group-theoretic moonshine theorem remains external; the celestial
framework supplies its conditional trace placement, not a global
\(\MtwFour\)-action on one K3 sigma model.
\end{corollary}

\begin{proof}
Compose Theorem~\ref{thm:uch-main}, Definition~\ref{def:cmb-combined-cel-algebra},
Theorem~\ref{thm:celestial-moonshine-bridge},
Theorem~\ref{thm:mathieu-celestial-correspondence}, and
Theorem~\ref{thm:shadow-tower-moonshine}. The three moonshine arrows
carry the same hypotheses as those theorems: the actual K3 symmetry
group in the geometric case, Gannon's virtual module for the full
\(\MtwFour\) character statement, and the maps
\(\mathcal R_n^{\mathrm{K3}}\) for any bar-depth/Rademacher comparison.
The conclusion is therefore exactly conditional trace placement.
\end{proof}

\begin{remark}[Companion status for Vol~III Theorem~\ref{thm:mathieu-moonshine}]
\label{rem:cmb-vol3-mathieu-upgrade}
Vol~III Theorem~\ref{thm:mathieu-moonshine} is tagged
\ClaimStatusProvedElsewhere\ (citing Eguchi--Ooguri--Tachikawa +
Gannon).  Corollary~\ref{cor:celestial-holography-recovers-moonshine}
does not upgrade that theorem to a second in-programme proof.  It gives
the conditional celestial placement:
\[
\text{Vol~III Theorem~\ref{thm:mathieu-moonshine}:}\quad
\ClaimStatusProvedElsewhere\ \text{(Gannon 2016)}
\;+\;
\ClaimStatusConditional\ \text{celestial trace placement
(Corollary~\ref{cor:celestial-holography-recovers-moonshine}).}
\]
The truth of the moonshine decomposition is still supplied by Gannon's
modular-form theorem.  The celestial bridge locates those virtual
\(\MtwFour\)-characters inside the K3 matter factor once the Gannon
module is supplied.
\end{remark}

\begin{remark}[Sources]
\label{rem:cmb-novelty}
\emph{Internal.} The conditional placement: celestial holography
applied to 4d twistor gravity coupled to K3 produces
\(\wtc\otimes \cA^{\mathcal{N}{=}4}_{K3}\); actual K3 sigma-model
symmetries act on the matter factor; Gannon's virtual \(\MtwFour\)
module supplies the full Mathieu twining characters; and the K3
polar/shadow data are recorded before any Rademacher/bar-depth
comparison is imposed.

\emph{External.}
The Mathieu moonshine coefficients \(A_n=2a_n\) are the EOT observation
proved by Gannon.  GHV classify actual K3 sigma-model symmetries; they
do not provide a global \(\MtwFour\)-action on one model.  The umbral
extension is constructed by Cheng--Duncan--Harvey and proved by
Duncan--Griffin--Ono.  The celestial \(\wtc\)-algebra of self-dual
gravity is supplied by Strominger, GHPS, BHS, and CPS.  The bridge
places the K3 characters as matter traces inside the celestial tensor
product.

\emph{Compared.}
The Vol~I ``Shadow = GW($\mathbb{C}^3$)'' row identifies the shadow tower
at $\kapch = \Psi$ with perturbative constant-map GW free energies on
$\CC^3$. The present chapter records the K3 polar/shadow comparison at
$\kapch = 2$ and isolates the additional datum needed to compare
bar-depth coefficients with moonshine Rademacher/Fourier coefficients.
These are distinct specializations of the shadow-tower machinery.
\end{remark}

% =====================================
\subsection*{Coda}
% =====================================

\begin{enumerate}
\item \emph{Statement.} The universal
celestial holography theorem applied to
$T_{\mathrm{sd grav}} \boxtimes \Sigma^{K3}$ produces
$\wtc \otimes \cA^{\mathcal{N}{=}4}_{K3}$. Actual K3 sigma-model
symmetries act trivially on the first factor and on the K3 matter
factor only. The full $\MtwFour$ twining functions enter through
Gannon's virtual Mathieu module, not through a geometric action on one
K3 sigma model. The polar and shadow data match the Mathieu mock
modular forms; the full Rademacher/bar-depth comparison requires the
maps \(\mathcal R_n^{\mathrm{K3}}\).

\item \emph{Scope.} The bridge produces a celestial chiral algebra and
a recovery of moonshine in the celestial sector; identification of the
celestial $\wtc$-algebra with a moonshine $\MtwFour$-module, equality
with any of the six Vol~III algebraizations
$G(K3 \times E)$, and statements about the Conway or Monster cases lie
outside its scope. The group-theoretic moonshine theorem itself
remains Gannon's.

\item \emph{Resolves.}
(a) The tensor-product placement of K3 moonshine inside the celestial
matter sector.
(b) The polar/shadow part of the aspirational shadow/$L$-function link
from Volume~I in the K3 lane.
(c) The Vol~III $\MtwFour$-action-on-$Y(\frakg_{K3})$-representations
conjecture (conj:m24-yangian-action) addresses the Yangian
representation category and is distinct from the celestial chiral
algebra construction here; both factor through the K3 Mukai lattice.

\item \emph{Consequence.} Granted the universal celestial
holography theorem, the Costello--Paquette K3 chiral algebra assignment,
and the Gannon virtual Mathieu module, the corresponding twined K3
characters appear as matter traces in the celestial sector.
\end{enumerate}

%=======================================================================
% End of celestial_moonshine_bridge.tex
%=======================================================================
